\documentclass[11pt]{article}

\usepackage[margin=25mm]{geometry}

\usepackage{amsmath,amssymb,fontspec,unicode-math}

\IfFontExistsTF{DejaVu Serif}{\setmainfont{DejaVu Serif}}{\setmainfont{Latin Modern Roman}}

\IfFontExistsTF{Latin Modern Math}{\setmathfont{Latin Modern Math}}{\setmathfont{DejaVu Math TeX Gyre}}

\usepackage{xurl,hyperref}

\setlength{\parindent}{0pt}\setlength{\parskip}{6pt}

\begin{document}

\title{Delving into Cascaded Instability: A Lipschitz Continuity View on Image Restoration and Object Detection Synergy}

\author{ELVIA preliminary AI analysis}\date{21 October 2026}\maketitle

Preliminary assessment: Minor.\par

Independent scientific review remains pending. Author adoption has not been requested.

The stability discussion confuses maximum directional gain with necessary amplification, and Remark 3 uses a training-point condition to control a global Jacobian maximum. Released training recipes also differ from the displayed objective.

\section*{Review scope}

Latest-version, claim-focused review of the cited author-paper passages and applicable current author-release paths. The stated counterexamples and scope controls were checked analytically; no author program, old verifier, training run or benchmark was executed.

\section*{Problem 1: Operator norms do not lower-bound every directional gain}

Source: Section 3.1, PDF p.4, paragraph before Remark 1

Assessment: supported; impact: minor

\subsection*{Source claim (paraphrase)}

Section 3.1 uses a restoration Jacobian norm above one to infer amplification of slight input changes, then invokes a larger detector norm.

\subsection*{Analysis}

Take \(A=\operatorname{diag}(2,1/10)\) and \(M=\operatorname{diag}(20,1/10)\). Their operator norms are \(2\) and \(20\), but the cascade maps \(\epsilon e_2\) to \((\epsilon/100)e_2\). These smooth linear maps therefore contract that finite perturbation despite large maximum gains.

\subsection*{Finding and scope}

Necessary amplification in every direction does not follow. Existence of amplified directions and the empirical observations remain possible.

\subsection*{Proposed correction}

The valid upper bound is \(\operatorname{Lip}(D\circ R)\leq\operatorname{Lip}(D)\operatorname{Lip}(R)\). A particular first-order gain is \(\|J_DJ_Rv\|/\|v\|\); a universal local lower bound requires a minimum-singular-value condition.

An explicit local correction is described; author adoption and a corresponding published erratum are not established.

\section*{Problem 2: Training-point alignment does not control the global envelope}

Source: Section 4.1, Remark 3, PDF p.6; Appendix A, PDF p.13

Assessment: supported; impact: minor

\subsection*{Source claim (paraphrase)}

Remark 3 assumes restoration-gradient alignment at a training point \(x^\star\), then bounds the derivative of the global Jacobian maximum \(C(\theta)\).

\subsection*{Analysis}

For \(B_\theta(x)=x+\theta(x^2-x)\) on \([0,1]\), \(0<\theta<1/2\), use losses \(L_r=2-\theta\), \(L_d=3-2\theta\), and weight \(1\). Then \(g_r=-1\), \(g_d=-2\), and at training point \(x^\star=0\) the required alignment is \(1\). Nevertheless \(C(\theta)=1+\theta\), attained at \(x=1\). Gradient flow has \(\dot\theta=3\), so \(\dot C=3\), while the printed bound is \(1\). Correcting only the detection-term sign does not repair the missing maximizer condition.

\subsection*{Finding and scope}

The stated sufficient condition does not establish this ancillary stability bound. It does not refute improved performance of the joint model.

\subsection*{Proposed correction}

Under the usual compactness and differentiable-envelope conditions, require alignment at every active maximizer. Use the signed detection contribution \[\max_{x\in\arg\max\phi}[-\langle\nabla_\theta\phi,g_d\rangle]\]. Strict decrease needs this term to be smaller than the restoration contribution.

An explicit local correction is described; author adoption and a corresponding published erratum are not established.

\section*{Problem 3: Public losses differ from the displayed parameter-Jacobian objective}

Source: Section 4.2, Remark 4 and full objective, PDF p.7; Table 4, PDF p.9; released loss path

Assessment: supported; impact: minor

\subsection*{Source claim (paraphrase)}

Section 4.2 uses \(L_d+\lambda L_r+\lambda_p\|\nabla_\theta f_\theta(x)\|\), with a Charbonnier restoration loss and stated weights \(\lambda=10\), \(\lambda_p=0.01\).

\subsection*{Analysis}

The public recipes assemble weighted absolute-error/perceptual reconstruction terms and clipping or loss-gradient penalties. A loss-gradient norm is not generally the model-output parameter-Jacobian norm in the equation; clipping is an optimizer operation rather than that objective. No configuration-to-table mapping establishes which historical experiment used these recipes.

\subsection*{Finding and scope}

The inspected recipe is insufficient to reproduce the exact displayed objective. No published table entry is numerically contradicted.

\subsection*{Proposed correction}

Provide the exact loss, norm, reduction and experiment configuration. An operator norm must not silently become a Frobenius estimator; implementing another regularizer requires its own interpretation and evaluation.

An explicit local correction is described; author adoption and a corresponding published erratum are not established.

\section*{Limits}

Preliminary AI assessment; independent human scientific adjudication remains pending.\par

An analytic counterexample or a release discrepancy does not by itself remeasure the paper’s reported experiments.\par

\section*{Source identity}

arXiv:2510.24232v3 (current author version; updated 2026-08-13T06:17:50Z)

\url{https://arxiv.org/pdf/2510.24232v3}

Reviewed PDF SHA-256: \nolinkurl{690c9485d07176a4a77f80eabe2c0a77636de60ce0f5f1593e6c379f34478604}

The selected author/publisher source is the same work as the original case, matched by exact title, author identity and official version channel. The exact selected PDF SHA-256 was recomputed in the latest-version pass for Issue \#27. The PDF is a retained author/publisher copy; its original retrieval time is stated only when supported by matching acquisition metadata. File modification time was not substituted.

\end{document}
